\newcommand{\DiamIonThree}{17}
\newcommand{\DiamIonThreePlusOne}{18}
\newcommand{\MeanIonThree}{10.72}
\newcommand{\DiamIbmThree}{19}
\newcommand{\MeanIbmThree}{12.34}
\newcommand{\BFSMinutes}{25}
\newcommand{\OddRelIon}{11}
\newcommand{\OddRelIonMinusOne}{10}
\newcommand{\OddRelIonHalf}{6}
\newcommand{\OddRelIonHalfMinusOne}{5}
\newcommand{\OddRelIbm}{9}
\newcommand{\TwoQThreeDist}{13{,}824/373{,}248/5{,}971{,}968/38{,}486{,}016/42{,}052{,}608/5{,}971{,}968/27{,}648 for $0/\dots/6$ entangling gates}
\newcommand{\TwoQThreeMean}{3.51}
\newcommand{\TwoQThreeMax}{6}
\newcommand{\OptIon}{10.4}
\newcommand{\FaithIon}{47.1}
\newcommand{\EWeightIon}{28.3}
\newcommand{\EWsixteenIon}{10.4}
\newcommand{\FaithIonSec}{51}
\newcommand{\EWeightIonSec}{0.54}
\newcommand{\EWsixteenIonSec}{0.24}
\newcommand{\GapFactorIon}{4.5}
\newcommand{\SpeedupIon}{210}
\newcommand{\WStarIon}{12}
\newcommand{\OptIbm}{12.2}
\newcommand{\FaithIbm}{36.8}
\newcommand{\EWeightIbm}{51.1}
\newcommand{\EWsixteenIbm}{12.2}
\newcommand{\FaithIbmSec}{17}
\newcommand{\EWeightIbmSec}{0.22}
\newcommand{\EWsixteenIbmSec}{0.09}
\newcommand{\GapFactorIbm}{3.0}
\newcommand{\SpeedupIbm}{190}
\newcommand{\WStarIbm}{16}
\newcommand{\GridEightFive}{46.1}
\newcommand{\GridEightInf}{17.9}
\newcommand{\MeanIonTwoEven}{5.90}
\newcommand{\MeanIonTwoOdd}{5.88}
\newcommand{\BetaTwo}{0.937}
\newcommand{\LMixTwo}{133}
\newcommand{\MixLThree}{100}
\newcommand{\LBFourIon}{11.1}
\newcommand{\LBFourIbm}{12.9}
\newcommand{\LBThreeCheck}{7.9}
\newcommand{\CWFourIon}{56.7}
\newcommand{\CWFourIonTwoQ}{27.6}
\newcommand{\CWFourSec}{5.6}
\newcommand{\EWFourSec}{0.8}
\newcommand{\FaithFourSec}{65}
\newcommand{\EWFourIon}{110.3}
\newcommand{\EWFourIonTwoQ}{44.2}
\newcommand{\FaithFourIon}{78.8}
\newcommand{\FaithFourIonTwoQ}{33.3}
\newcommand{\QFourIon}{164.9}
\newcommand{\QFourIonTwoQ}{47.3}
\newcommand{\FourImprovementIon}{49\%}
\newcommand{\QThreeIon}{127.2}
\newcommand{\CWFourIbm}{59.1}
\newcommand{\CWFourIbmTwoQ}{27.1}
\newcommand{\EWFourIbm}{130.9}
\newcommand{\EWFourIbmTwoQ}{47.1}
\newcommand{\FaithFourIbm}{76.0}
\newcommand{\FaithFourIbmTwoQ}{32.4}
\newcommand{\QFourIbm}{140.9}
\newcommand{\QFourIbmTwoQ}{47.3}
\newcommand{\FourImprovementIbm}{45\%}
\newcommand{\QThreeIbm}{104.0}
\newcommand{\CWBqIonN}{30}
\newcommand{\CWBqIonTwoQ}{27.0}
\newcommand{\CWBqIonFewer}{16}
\newcommand{\CWBqIonMore}{13}
\newcommand{\CWBqIonTotalRatio}{0.46}
\newcommand{\CWBqIbmN}{30}
\newcommand{\CWBqIbmTwoQ}{28.1}
\newcommand{\CWBqIbmFewer}{25}
\newcommand{\CWBqIbmMore}{4}
\newcommand{\CWBqIbmTotalRatio}{0.37}
\newcommand{\GapFindings}{\paragraph{Findings.} (1) On these benchmarks the original method finishes \GapFactorIon$\times$ (ion trap) and
\GapFactorIbm$\times$ (IBM-style) above the exact optimum. (2) EW-$w$ reaches the
\emph{proven} optimum on all 20 circuits once $w\ge\WStarIon$ (ion trap) or $w\ge\WStarIbm$ (IBM-style), in a fraction of a second.
This is the cascade of Prop.~\ref{prop:cascade}. The transition is sharp, between $w=8$ and $w=12$ (ion trap) or $w=16$ (IBM-style). (3) The grid (Fig.~\ref{fig:grid})
shows that the \emph{database radius}, not the window or the search strategy, is the binding constraint. At $r=5$, enlarging the
window from 8 to 16 lowers the mean excess only from 46 to 43 gates. At $w=8$, raising $r$ from 5 to $\infty$ lowers it from 46 to 18,
and for $w\ge12$ with $r\ge8$ the excess is at most 2 gates (ion trap, and from 72 to $\le$\,8 for the IBM-style set). Theorem~\ref{thm:exact} predicts
both requirements, windows long enough \emph{and} a database deep enough. The original method is limited by the second.
}
\newcommand{\FourFindings}{\paragraph{Findings.} (1) \emph{Convexity matters more than anything else at $n\ge4$.} With the same exact 3-qubit
oracle, contiguous windows (EW-16) give \EWFourIon\ (ion trap) and \EWFourIbm\ (IBM-style) gates, no better than the
published method. With DAG-convex windows (CW-24) the same oracle gives \CWFourIon\ and \CWFourIbm\ gates, i.e.\
\FourImprovementIon\ and \FourImprovementIbm\ fewer than the published results, in about \CWFourSec\,s per circuit. (2)
Two-qubit counts drop from 43 (published) to \CWFourIonTwoQ\ RXX and \CWFourIbmTwoQ\ CZ, although CW optimises
total length, not two-qubit count. On the first \CWBqIonN\ circuits, where we also ran BQSKit level~4
(Table~\ref{tab:four}), CW-24 averages \CWBqIonTwoQ\ RXX against BQSKit's \BqFourIonTwoQ\ (fewer on \CWBqIonFewer, more
on \CWBqIonMore) and \CWBqIbmTwoQ\ CZ against \BqFourIbmTwoQ\ (fewer on \CWBqIbmFewer, more on \CWBqIbmMore), with
$\times$\CWBqIonTotalRatio\ and $\times$\CWBqIbmTotalRatio\ of BQSKit's total gates and exact outputs. The journal
reports 28 and 37 two-qubit gates for BQSKit. (3) The certified lower bounds
(\LBFourIon, \LBFourIbm) leave a gap of at most a factor of about 4 (\NjpIonCliff/\LBFourIon\ and \NjpIbmCliff/\LBFourIbm). Closing it, either by tightening the bound (deeper spheres,
meet-in-the-middle) or by better upper bounds (4-qubit oracles via symmetry reduction), is an open problem that our framework
makes precise.
}
